% !TeX root = ../thesis.tex
% ************************** Abstract *****************************
% Use the `abstract' option in the document class to print only the title page and abstract.
\begin{abstract}

Vietnamese Visual Question Answering requires a model that both understands the
image and reads the Vietnamese text in it, including its diacritics.
Open-source vision-language models (VLMs), however, are trained mostly on English,
and fully fine-tuning such a model is beyond the resources of a student team. This
thesis studies how to adapt a VLM to Vietnamese using only parameter-efficient
fine-tuning on a single Google Colab GPU.

We select Qwen2.5-VL-3B-Instruct as the baseline, reproduce its published result on
the TextVQA validation set (VQA accuracy 79.60 vs.\ 79.3 reported), and then evaluate
it zero-shot on two Vietnamese datasets: ViVQA (everyday images) and ViTextVQA (images containing text).
Zero-shot, the model reaches only 22.09 and 29.12 Exact Match (EM), mostly because its
answers have the wrong format or are in the wrong language. Fine-tuning with 4-bit
QLoRA, which trains only 0.196\% of the parameters, raises EM to 76.64 on ViVQA and
50.26 on ViTextVQA.

The proposed improvement is \textbf{OCR-aware QLoRA}: the text in the image is read
with PaddleOCR and VietOCR, sorted in reading order and inserted into the prompt
during training as well as at inference. On the full ViTextVQA test set of 10{,}028 questions, the
method reaches 53.35 EM and 73.89 ANLS, 3.09 EM points above plain QLoRA on the same
data (McNemar test $p = 6.7\times10^{-21}$). The gain is stable across three seeds,
and the result is well above the best previously published one (25.48 EM). A $2\times2$ factorial
experiment shows that the gain appears only when OCR is present both in training and
at inference (interaction $+4.9$ EM): the contribution lies in teaching the model to
use OCR, not in adding OCR to the prompt.

The thesis includes ablations, an automatic error analysis checked against a blind
manual review, and a hallucination measurement with POPE-vi, a set of 4{,}985
Vietnamese questions that we built. The results show that EM underestimates the
model (about 55\% of EM errors are semantically correct), that Vietnamese text
reading errors still cost about 17 EM points, and that hallucination accounts for
only about 2--4\% of the answers. Every run uses a single L4 GPU with at most 8.4~GB
of VRAM and is logged in a reproducibility package. Finally, the system is packaged
as a web prototype that shows the OCR text as evidence for each answer.

\vspace{1em}
\noindent\textbf{Keywords:} Vietnamese VQA, vision-language model, QLoRA,
parameter-efficient fine-tuning, OCR, scene text, hallucination.

\end{abstract}
